\appendix

\chapter{Chứng minh Mệnh đề~\ref{prop:decomposition}}
\label{appendix:taylor_proof}

Chương này chứng minh chi tiết Mệnh đề~\ref{prop:decomposition} (Mục~\ref{sec:definition_AG}): dưới chế độ thích ứng nhanh (Giả thuyết~\ref{asm:fast_adaptation}), $\Delta M = -C_b + O\big((T\alpha)^2\big)$.

\section{Thiết lập và khai triển bậc hai}

Ký hiệu $\theta_0 = (\theta_0^b, \theta_0^h)$. Xét hàm mất mát kỳ vọng trên tập truy vấn $L(\phi) := \mathbb{E}_{Q_i\sim\mathcal{T}_i}[\mathcal{L}(Q_i,\phi)]$, giả sử khả vi liên tục hai lần trong một lân cận mở chứa $\theta_0$ và toàn bộ quỹ đạo. Đặt $_b := \nabla_b L(\theta_0)$, $g_h := \nabla_h L(\theta_0)$, $H_{bb} := \nabla^2_{bb} L(\theta_0)$, $H_{bh} := \nabla^2_{bh} L(\theta_0)$, $H_{hh} := \nabla^2_{hh} L(\theta_0)$.

Gọi $\Delta b_T := \phi_T^b - \theta_0^b$, $\Delta h_T := \phi_T^h - \theta_0^h$ (nhánh thích ứng toàn phần) và $\Delta h_T^f := \phi_T^{f,h} - \theta_0^h$ (nhánh can thiệp). Vì backbone bị đóng băng ở nhánh can thiệp, $\Delta b_T^f = 0$.

Khai triển Taylor bậc hai của $L$ quanh $\theta_0$ cho mỗi nhánh:
\begin{align}
    L(\phi_T) &= L(\theta_0) + g_b^\top \Delta b_T + g_h^\top \Delta h_T \nonumber\\
    &\quad + \tfrac{1}{2}\Delta b_T^\top H_{bb} \Delta b_T + \Delta b_T^\top H_{bh} \Delta h_T + \tfrac{1}{2}\Delta h_T^\top H_{hh} \Delta h_T + o(\|\Delta\|^2), \label{eq:taylor_full}\\[4pt]
    L(\phi_T^f) &= L(\theta_0) + g_h^\top \Delta h_T^f + \tfrac{1}{2}{\Delta h_T^f}^\top H_{hh} \Delta h_T^f + o(\|\Delta h_T^f\|^2). \label{eq:taylor_frozen}
\end{align}
Trừ~\eqref{eq:taylor_frozen} cho~\eqref{eq:taylor_full} và dùng định nghĩa $\Delta M = L(\phi_T^f) - L(\phi_T)$ (Phương trình~\eqref{eq:adaptation_gain}, đã lấy kỳ vọng theo $Q_i$), số hạng $L(\theta_0)$ triệt tiêu và ta thu được chính xác Phương trình~\eqref{eq:decomposition}:
\begin{equation}
    \begin{split}
        \Delta M = & -\,\underbrace{\Big[g_b^\top \Delta b_T + \tfrac{1}{2}\Delta b_T^\top H_{bb} \Delta b_T + \Delta b_T^\top H_{bh} \Delta h_T\Big]}_{C_b} \\
    &+ \underbrace{\Big[g_h^\top(\Delta h_T^f - \Delta h_T) + \tfrac{1}{2}\big({\Delta h_T^f}^\top H_{hh} \Delta h_T^f - \Delta h_T^\top H_{hh} \Delta h_T\big)\Big]}_{R_h},
    \end{split}
    \label{eq:decomposition}
\end{equation}

\section{Bậc độ lớn của các số hạng dịch chuyển}

\begin{lemma}[Dịch chuyển tích lũy $O(T\alpha)$]
\label{lem:displacement_order}
Giả sử $\|\nabla_\phi \mathcal{L}(\phi,S_i)\| \le G$ trong lân cận quỹ đạo. Khi đó $\|\Delta b_T\|, \|\Delta h_T\|, \|\Delta h_T^f\| = O(T\alpha)$.
\end{lemma}
\begin{proof}
Mỗi bước cập nhật inner-loop~\eqref{eq:inner_loop} dịch chuyển tham số một lượng $\alpha\|\nabla_\phi\mathcal{L}\| \le \alpha G$. Sau $T$ bước, theo bất đẳng thức tam giác, tổng dịch chuyển tích lũy bị chặn bởi $T\alpha G = O(T\alpha)$, áp dụng như nhau cho cả ba đại lượng vì cả hai nhánh đều thực hiện đúng $T$ bước cập nhật dạng này (nhánh can thiệp chỉ khác ở việc phần backbone bị ghi đè về $\theta_0^b$ sau mỗi bước, không ảnh hưởng đến bậc của phần head).
\end{proof}

\begin{lemma}[Độ lệch giữa hai quỹ đạo head là $O((T\alpha)^2)$]
\label{lem:head_divergence}
Đặt $\delta_j := \Delta h_j - \Delta h_j^f$. Dưới điều kiện $\nabla_h\mathcal{L}(\cdot,S_i)$ Lipschitz với hằng số $L$ theo $\phi$ (hệ quả của Giả thuyết~\ref{asm:adjoint_stability}(ii)), ta có $\delta_T = O\big((T\alpha)^2\big)$. Đặc biệt, tại $T=1$, $\delta_1 = 0$ \textbf{tuyệt đối}.
\end{lemma}
\begin{proof}
Tại $j=1$: cả hai nhánh cùng xuất phát từ $\phi_0 = \phi_0^f = \theta_0$, nên bước gradient đầu tiên của head trùng nhau tuyệt đối, $\delta_1 = 0$.

Với $j \ge 2$, từ quy tắc cập nhật:
\[
    \delta_j = \delta_{j-1} - \alpha\Big[\nabla_h\mathcal{L}(\phi_{j-1},S_i) - \nabla_h\mathcal{L}(\phi_{j-1}^f,S_i)\Big].
\]
Hai đối số $\phi_{j-1} = (\phi_{j-1}^b,\phi_{j-1}^h)$ và $\phi_{j-1}^f=(\theta_0^b,\phi_{j-1}^{f,h})$ khác nhau ở backbone một lượng $\Delta b_{j-1} = O((j-1)\alpha)$ (Bổ đề~\ref{lem:displacement_order}) và ở head một lượng $\delta_{j-1}$. Do tính Lipschitz,
\begin{align*}
    \big\|\nabla_h\mathcal{L}(\phi_{j-1},S_i) - \nabla_h\mathcal{L}(\phi_{j-1}^f,S_i)\big\| &\le L\big(\|\Delta b_{j-1}\| + \|\delta_{j-1}\|\big)\\
                                                                                            &= O\big((j-1)\alpha + \|\delta_{j-1}\|\big).
\end{align*}
Suy ra
\[
    \|\delta_j\| \le (1+\alpha L)\|\delta_{j-1}\| + O\big((j-1)\alpha^2\big).
\]
Với $\delta_1=0$ và khai triển đệ quy, do $T\alpha\ll1$ (Giả thuyết~\ref{asm:fast_adaptation}) nên $(1+\alpha L)^{T} = 1+O(T\alpha)$ vẫn bị chặn:
\begin{align*}
    \|\delta_T\| &\le \sum_{k=1}^{T-1} O(k\alpha^2)\cdot\big(1+O(\alpha)\big)^{T-1-k}\\
    &= O(\alpha^2)\sum_{k=1}^{T-1} k \cdot\big(1+O(T\alpha)\big)\\
    &= O(\alpha^2 T^2) = O\big((T\alpha)^2\big).
\end{align*}
\end{proof}

\section{Bậc độ lớn của $C_b$ và $R_h$, và kết luận}

Từ Bổ đề~\ref{lem:displacement_order}: số hạng tuyến tính $g_b^\top\Delta b_T = O(T\alpha)$ là bậc chủ đạo trong $C_b$; hai số hạng bậc hai còn lại của $C_b$ đều là $O((T\alpha)^2)$. Vậy $C_b = O(T\alpha)$.

Từ Bổ đề~\ref{lem:head_divergence}: số hạng tuyến tính của $R_h$ là $g_h^\top(\Delta h_T^f - \Delta h_T) = -g_h^\top\delta_T = O\big((T\alpha)^2\big)$; các số hạng bậc hai còn lại của $R_h$ (hiệu hai dạng toàn phương) cũng có bậc $O\big((T\alpha)^2\big)$ vì $\delta_T$ là bậc chủ đạo của hiệu $\Delta h_T^f$ và $\Delta h_T$. Vậy $R_h = O\big((T\alpha)^2\big)$, và
\[
    \Delta M = -C_b + R_h = -C_b + O\big((T\alpha)^2\big),
\]
đúng như Phương trình~\eqref{eq:leading_order}, hoàn tất chứng minh Mệnh đề~\ref{prop:decomposition}. $\blacksquare$

Đặc biệt, với $T=1$, Bổ đề~\ref{lem:head_divergence} cho $\delta_1=0$ tuyệt đối, kéo theo $R_h=0$ \emph{chính xác} (không chỉ là vô cùng bé), khớp với nhận xét tại Mục~\ref{sec:definition_AG}.
